\section{Discussion}

% Kritische Diskussion

% -----

% • Wozu macht man Fits gemessener Daten?
% • Wie lautet der zentrale Grenzwertsatz? Am besten in einem Statistik-Lehrbuch nachschlagen!
% • Bei einer Reihe von N Messungen in einem Messpunkt mit Standardabweichung σ, wie groß
% ist der statistische Fehler δ des Mittelwerts?
% • Bei mehreren Reihen mit je N Messungen im selben Messpunkt, wie ist die Verteilung der
% Mittelwerte?
% • Was zeichnet die ”
% Least Squares“-Methode aus? Wie ist χ
% 2 definiert?
% • Wie schätzt man quantitativ die Qualität eines Fits ab? Was für Schlüsse lassen sich aus der Abschatzung ziehen?

\subsection{Experiment A}

In the first experiment part, we were able to find plausible Lorentz curves for the first two harmonics.
The curves along with the data can be seen in \autoref{fig:analysis_exp_a}.
As can be seen in the graph of the data, it was not possible to properly measure the third resonance peak.
The measured signal contains too much noise to see the peak.
This also leads to a failed fit, and we simply see a flat \quoteapprox{Lorentz curve}.

While the curves for the first two resonances look alright visually, goodness of fit analysis produces very high $\chi^2_\mathrm{red}$ values.
This statistic measures how far away the measurements are from their expected positions.
Derived from this is the $p$-value of the fit, which represents the probability of measuring data at least this extreme (far away from the expectation), given that the model we found is correct.

To the numerical precision of the calculation software, these values are zero, or end up very small (see \autoref{tab:analysis_stats_a}).
Using the $p$-values and any reasonable significance level (such as \num{0.05}), we must conclude that the model is not correct, and that in all three cases the measurements did not come from the distributions described by the parameters we calculated in \autoref{tab:analysis_params_a}.

As was noted in the analysis, visually inspecting the residuals of the fit also gives an idea about any systematic over- or underestimates of the data in certain regions.
In our case, the fit functions for the first two harmonics do produce patterns in their residuals.
Inspecting \autoref{fig:analysis_exp_a_residuals}, we see that the model of the first harmonic overestimates the size of the tails and underestimates towards the center/peak of the curve.
The opposite happens with the second fit function around \qty{825}{\hertz}, where the model underestimates the tails and center but overestimates in-between.

To conclude this experiment, we attempted to calculate the quality of the resonator ($Q$-factor) for each peak (even though the fit parameters are likely not accurate).
Because the fit parameters are most likely not well determined, we do not expect the theoretical $Q$-factors to be correct.
As we see in \autoref{tab:analysis_Q_a}, the factors determined from the fit parameters are illogical and do not fit with the experimentally determined ones.
The large errors in $Q_\mathrm{fit}$ mean the $\sigma$-deviations are very low, but viewing both the relative and absolute deviation together it is easy to conclude that the calculations were not successful.

\subsection{Experiment B}

In part B we analyzed ten resonance curves which were measured at different temperatures. For each of these we performed fits using Lorentz curves, like in experiment A.
These fits were used to determine the exact positions of the peaks, which could be graphed against the temperature of the read during each of the spectrum measurements.
The resulting graph is shown in \autoref{fig:analysis_exp_b_temp_graph}.
A linear fit was performed on this data to obtain the fit parameters:
\begin{gather}
    \boxed{
        \begin{aligned}
             & m &  & = -0.001441 ± 0.000014 \\
             & b &  & = 125.0490 ± 0.0009
        \end{aligned}
    }
\end{gather}

For this relationship, the slope is called the temperature coefficient.
It describes how an increase in temperature affects the frequency of the resonance maximum.
We found this to have a negative sign, meaning an increase in temperature decreases the frequency at which the resonance reaches its peak.

